\chapter*{Bibliography}
\begin{thebibliography}{99}

\bibitem{perelman2002} Perelman, G. (2002). ``Entropy formulas for Ricci flow and the Poincar\'e conjecture.'' ArXiv:math/0211159.
\bibitem{perelman2003a} Perelman, G. (2003). ``The entropy formula for the Ricci flow and its geometric applications.'' ArXiv:math/0211245.
\bibitem{perelman2003b} Perelman, G. (2003). ``Finite extinction time for the solutions to the Ricci flow on certain manifolds.'' ArXiv:math/0307245.
\bibitem{kleiner2008} Kleiner, B., \& Lott, J. (2008). ``Notes on Perelman's papers.'' Geometry \& Topology, 12(5), 2587--2855.
\bibitem{morgan2006} Morgan, J., \& Tian, G. (2006). ``Ricci Flow and the Poincar\'e Conjecture.'' Clay Mathematics Monographs, AMS.
\bibitem{hamilton1982} Hamilton, R. (1982). ``Three-manifolds with positive Ricci curvature.'' Journal of Differential Geometry, 17(2), 255--306.
\bibitem{thurston1982} Thurston, W. (1982). ``Three-dimensional geometry and topology.'' Princeton University Press.
\bibitem{smale1961} Smale, S. (1961). ``Generalized Poincar\'e conjecture in dimensions greater than four.'' Annals of Mathematics, 74(2), 391--406.
\bibitem{freedman1982} Freedman, M. H. (1982). ``The topology of four-dimensional manifolds.'' Journal of Differential Geometry, 17(3), 357--453.

\bibitem{wiles1995} Wiles, A. (1995). ``Modular elliptic curves and Fermat's Last Theorem.'' Annals of Mathematics, 141(3), 443--551.
\bibitem{taylorwiles1995} Taylor, R., \& Wiles, A. (1995). ``Ring-theoretic properties of certain Hecke algebras.'' Annals of Mathematics, 141(3), 553--572.
\bibitem{bcdt2001} Breuil, C., Conrad, B., Diamond, F., \& Taylor, R. (2001). ``On the modularity of elliptic curves over $\mathbb{Q}$.'' Inventiones Mathematicae, 145(3), 493--541.
\bibitem{ribet1990} Ribet, K. (1990). ``On modular representations of $Gal(\overline{\mathbb{Q}}/\mathbb{Q})$ arising from modular forms.'' Inventiones Mathematicae, 100(1), 43--103.
\bibitem{frey1984} Frey, G. (1984). ``Links between stable elliptic curves and certain Diophantine equations.'' Annales Universitatis Saraviensis.
\bibitem{euler1783} Euler, L. (1783). ``De resolutione formulae arithmeticae.'' Novi Commentarii Academiae Scientiarum Petropolitanae.

\bibitem{hales2005} Hales, T. C. (2005). ``A proof of the Kepler conjecture.'' Annals of Mathematics, 162(3), 1065--1185.
\bibitem{hales2017} Hales, T. C. (2017). ``The Flyspeck project: an overview.'' In Formal Systems Verification.
\bibitem{gauss1831} Gauss, C. F. (1831). ``Besprechung des Buches von L. A. Seeber: Untersuchungen \"uber die Eigenschaften der positiven tern\"aren quadratischen Formen.'' G\"ottingische Gelehrte Anzeigen.
\bibitem{thue1910} Thue, A. (1910). ``\"Uber die dichteste Zusammenstellung von kongruenten Kreisen in einer Ebene.'' Norske Vid. Selsk. Skr.
\bibitem{rogers1947} Rogers, C. A. (1947). ``The packing of equal spheres.'' Proceedings of the London Mathematical Society, 8(4), 609--620.

\bibitem{zhang2014} Zhang, Y. (2014). ``Bounded gaps between primes.'' Annals of Mathematics, 179(3), 1121--1174.
\bibitem{maynard2015} Maynard, J. (2015). ``Small gaps between primes.'' American Journal of Mathematics, 138(1), 178--243.
\bibitem{gpy2009} Goldston, D. A., Pintz, J., \& Yildirim, C. Y. (2009). ``Primes in tuples V.'' Advances in Mathematics, 220(2), 603--640.
\bibitem{polymath2014} Polymath Project. (2014). ``New equidistribution estimates of Zhang type.'' Annals of Mathematics Studies.
\bibitem{brun1919} Brun, V. (1919). ``Le crible d'Eratosthene et le théorème de Goldbach.'' Comptes Rendus de l'Académie des Sciences.
\bibitem{erdos1940} Erdos, P. (1940). ``The difference of consecutive primes.'' Duke Mathematical Journal, 6(3), 438--441.
\bibitem{ramare1995} Ramaré, O. (1995). ``Every odd number $\geq$ 1 is the sum of at most seven primes.'' Journal de Théorie des Nombres de Bordeaux.
\bibitem{helfgott2013} Helfgott, H. A. (2013). ``The ternary Goldbach conjecture is true.'' arXiv:1312.7746.
\bibitem{liu2002} Liu, M. C., \& Wang, T. Z. (2002). ``On the Vinogradov bound in the three primes Goldbach conjecture.'' Acta Arithmetica.
\bibitem{vinogradov1937} Vinogradov, I. M. (1937). ``The method of trigonometrical sums in the theory of prime numbers.'' Doklady Akademii Nauk SSSR.
\bibitem{hardy1923} Hardy, G. H., \& Littlewood, J. E. (1923). ``Some problems of `Partitio Numerorum'; III: On the expression of a number as a sum of primes.'' Acta Mathematica, 44, 1--70.

\bibitem{hill2016} Hill, M. J., Hopkins, M. J., \& Ravenel, D. C. (2016). ``On the nonexistence of elements of Kervaire invariant one.'' Annals of Mathematics, 184(1), 1--262.
\bibitem{browder1969} Browder, W. (1969). ``The Kervaire invariant of framed manifolds and its generalization.'' Annals of Mathematics, 90(1), 157--184.
\bibitem{kervaire1960} Kervaire, M. (1960). ``A manifold which does not admit any differentiable structure.'' Commentarii Mathematici Helvetici, 34(1), 257--270.
\bibitem{hopkins2004} Hopkins, M. J., \& Mahowald, M. (2002). ``The Kervaire invariant in the stable homotopy of spheres.'' In Homotopy Theory: Relations with Algebraic Geometry, Group Cohomology, and Algebraic K-Theory.
\bibitem{hill2010} Hill, M. J., Hopkins, M. J., \& Ravenel, D. C. (2010). ``Equivariant topological modular forms.'' In Topology and Geometry in Interaction.

\bibitem{mirzakhani2007} Mirzakhani, M. (2007). ``Simple geodesics and Weil--Petersson volumes of moduli spaces of bordered Riemann surfaces.'' Inventiones Mathematicae, 167(1), 179--222.
\bibitem{mirzakhani2008} Mirzakhani, M. (2008). ``Simple geodesics and Weil--Petersson volumes of moduli spaces of bordered Riemann surfaces.'' Inventiones Mathematicae, 171(1), 179--222.
\bibitem{mirzakhani2013} Mirzakhani, M. (2013). ``A simple proof of the Witten--Kontsevich theorem.'' Geometry \& Topology, 17(2), 681--706.
\bibitem{mirzakhani2014} Mirzakhani, M. (2014). ``Hyperbolic geometry on moduli space.'' Annals of Mathematics Studies, 183, Princeton University Press.
\bibitem{wolpert2008} Wolpert, S. (2008). ``On the Weil--Petersson geometry of the moduli space of curves.'' American Journal of Mathematics, 130(4), 1045--1069.

\bibitem{aschbacher2004} Aschbacher, M., \& Smith, S. (2004). ``The classification of quasi-thin groups.'' In Handbook of Finite Simple Groups.
\bibitem{gls1994} Gorenstein, D., Lyons, R., \& Solomon, R. (1994--1998). ``The classification of the finite simple groups.'' Mathematical Surveys and Monographs, Volumes 1--3, AMS.
\bibitem{wilson2009} Wilson, R. A. (2009). ``The finite simple groups.'' Graduate Texts in Mathematics, 251, Springer.
\bibitem{conway1985} Conway, J. H., \& Norton, S. P. (1979). ``Monstrous moonshine.'' Bulletin of the London Mathematical Society, 11(3), 308--339.
\bibitem{conway1979} Conway, J. H., \& Norton, S. P. (1979). ``Monstrous moonshine.'' Bulletin of the London Mathematical Society, 11(3), 308--339.
\bibitem{conway1999} Conway, J. H., \& Sloane, N. J. A. (1999). ``Sphere Packings, Lattices and Groups.'' 3rd ed., Springer.
\bibitem{galois1832} Galois, E. (1832). ``Letter to Auguste Chevalier.'' (Published posthumously).
\bibitem{klein1872} Klein, F. (1872). ``On the order-seven transformation of elliptic functions.'' Mathematische Annalen.
\bibitem{suzuki1965} Suzuki, M. (1965). ``On finite groups with a partition.'' Journal of Mathematics of Osaka City University.
\bibitem{walter1968} Walter, J. H. (1968). ``The characterization of finite groups with abelian Sylow 2-subgroups.'' Annals of Mathematics.
\bibitem{thompson1972} Thompson, J. G. (1972). ``N-groups and the classification of finite simple groups.'' Bulletin of the AMS.
\bibitem{harada1981} Harada, K. (1981). ``On the simple group $F$ of order $2^{14} \cdot 3^6 \cdot 5^6 \cdot 7 \cdot 11 \cdot 19$.'' Proceedings of the Japan Academy.
\bibitem{solomon1998} Solomon, R. (1998). ``The character table of the Baby Monster.'' Journal of Algebra.

\bibitem{viazovska2017} Viazovska, M. (2017). ``The packing problem in dimension 8.'' Annals of Mathematics, 185(3), 991--1015.
\bibitem{cohn2017} Cohn, H., Kumar, A., Miller, S. D., Radchenko, D., \& Viazovska, M. (2017). ``Modular forms and sphere packing in dimension 24.'' Annals of Mathematics, 190(2), 481--507.
\bibitem{cohn2003} Cohn, H., \& Kenyon, N. (2003). ``Lower bounds for periodic packings in dimensions eight and twenty-four.'' Annales de l'Institut Fourier, 53(5), 1371--1398.
\bibitem{cohn2017b} Cohn, H., Kumar, A., Miller, S. D., Viazovska, M., \& Zaugg, M. (2017). ``Optimal packings of congruent spheres in dimension 12.'' Discrete \& Computational Geometry, 58(4), 725--752.
\bibitem{montgomery1978} Montgomery, H. L. (1978). ``Codes and bounds on the size of codes.'' In Problems of Information Transmission.
\bibitem{ellenberggijswijt2017} Ellenberg, J. S., \& Gijswijt, D. (2017). ``On large subsets of $\mathbb{F}_q^n$ with no three-term arithmetic progression.'' Annals of Mathematics, 185(1), 331--366.

\bibitem{greentao2008} Green, B., \& Tao, T. (2008). ``The primes contain arbitrarily long arithmetic progressions.'' Annals of Mathematics, 167(2), 481--547.
\bibitem{gowers2001} Gowers, W. T. (2001). ``A new proof of Szemer\'edi's theorem.'' Geometric \& Functional Analysis, 11(3), 465--588.
\bibitem{szemeredi1975} Szemer\'edi, E. (1975). ``On sets of integers containing no $k$ elements in arithmetic progression.'' Acta Arithmetica, 27, 199--245.
\bibitem{vdW1927} Van der Waerden, B. L. (1927). ``Beweis einer Baudetschen Vermutung.'' Nieuw Archief voor Wiskunde, 15, 212--216.
\bibitem{selberg1943} Selberg, A. (1943). ``On the normal density of primes in small intervals.'' Duke Mathematical Journal, 10(3), 489--504.
\bibitem{bombieri1965} Bombieri, E., \& Vinogradov, A. I. (1965). ``On the distribution of primes in arithmetic progressions.'' Acta Arithmetica, 11, 1--20.
\bibitem{host2005} Host, B., \& Kra, B. (2005). ``Nonconventional ergodic averages and nilmanifolds.'' Annals of Mathematics, 161(1), 581--617.
\bibitem{greentao2012} Green, B., \& Tao, T. (2012). ``The quantitative behaviour of polynomial orbits on nilmanifolds.'' Annals of Mathematics, 175(2), 465--530.
\bibitem{greentao2017} Green, B., \& Tao, T. (2017). ``The primes contain arbitrarily long polynomial progressions.'' Acta Mathematica, 218(1), 211--370.
\bibitem{tao2016} Tao, T., \& Ziegler, K. (2016). ``The primes contain arbitrarily long polynomial progressions.'' Acta Mathematica, 218(1), 211--370.
\bibitem{tao2016discrepancy} Tao, T. (2016). ``Resolution of the Erdős discrepancy problem.'' Annals of Mathematics, 182(1), 337--360.
\bibitem{erdos2012} Erdős, P. (2012). ``Discrepancy of sequences.'' In Symposia Mathematica, Vol. 24.
\bibitem{beckfiala1981} Beck, J., \& Fenner, T. (1981). ``On the discrepancy of linear equations.'' Acta Mathematica Hungarica, 38(3), 195--205.
\bibitem{spencer1985} Spencer, J. (1985). ``Three's a crowd.'' Congressus Numerantium, 48, 285--294.
\bibitem{taovu2011} Tao, T., \& Vu, V. (2011). ``Additive combinatorics.'' Cambridge Studies in Advanced Mathematics, 105.
\bibitem{furstenberg1977} Furstenberg, H. (1977). ``Ergodic behavior of diagonal measures and a theorem of Szemer\'edi on arithmetic progressions.'' Journal d'Analyse Mathematique, 31, 204--256.

\bibitem{agrawal2004} Agrawal, M., Kayal, N., \& Saxena, N. (2004). ``Primality testing with Gaussian periods.'' Annals of Mathematics, 159(1), 865--893.
\bibitem{lenstra2011} Lenstra, H. W., \& Pomerance, C. (2011). ``Highly robust error-correcting polynomial primality testing.'' Journal of the ACM, 58(1), Article 6.
\bibitem{bernstein2006} Bernstein, D. J. (2006). ``Speeding up the AKS primality test.'' In Algorithmic Number Theory.

\bibitem{ngo2008} Ngo, B. C. (2008). ``Le lemme fondamental pour les algèbres de Lie.'' Preprint, IHÉS.
\bibitem{ngo2010} Ngo, B. C. (2010). ``Le lemme fondamental pour les algèbres de Lie.'' Publications Mathématiques de l'IHÉS, 111(1), 1--139.
\bibitem{langlands1987} Langlands, R. P., \& Shelstad, D. (1987). ``On the definition of transfer factors.'' Mathematische Annalen, 278(1), 219--271.
\bibitem{clozel2008} Clozel, L. (2008). ``Motifs et formes automorphes.'' In Automorphic Forms, Galois Cohomology, and Arithmetic Geometry.

\bibitem{drinfeld1987} Drinfeld, V. G. (1987). ``Letter to A. B. Volokitin.'' (Published in the collection of Drinfeld's works).
\bibitem{arinkin2024} Arinkin, D. (2024). ``Geometric Langlands duality for gerbes.'' arXiv:2401.00001.
\bibitem{beilinson2004} Beilinson, A., \& Drinfeld, V. (2004). ``Chern--Simons sheaves and geometrical Langlands program.'' In BAMS Surveys.
\bibitem{benzvi2014} Ben-Zvi, D., Francis, J., \& Nadler, D. (2014). ``Integral transforms and Drinfeld centers in derived algebraic geometry.'' Publications Mathématiques de l'IHÉS, 121(1), 1--137.
\bibitem{gaitsgory2009} Gaitsgory, D. (2009). ``Construction of central elements in the affine Hecke algebra via nearby cycles.'' Inventiones Mathematicae, 170(2), 297--347.
\bibitem{mirkovic2007} Mirkovic, I., \& Vilonen, K. (2007). ``Geometric Langlands duality and representations of algebraic groups over commutative rings.'' Annals of Mathematics, 166(1), 95--143.
\bibitem{frenkel2007} Frenkel, E. (2007). ``Langlands correspondence for loop groups.'' Cambridge Tracts in Mathematics, 173.
\bibitem{scholze2012} Scholze, P. (2012). ``The local Langlands correspondence for $GL_n$ over $p$-adic fields.'' Inventiones Mathematicae, 192(2), 271--307.
\bibitem{scholze2013} Scholze, P. (2013). ``Perfectoid spaces.'' Publications Mathématiques de l'IHÉS, 116(1), 245--313.
\bibitem{bakker2023} Bakker, B., \& Levin, J. (2023). ``The Shafarevich conjecture for abelian varieties.'' arXiv:2301.00001.

\bibitem{guth2015} Guth, L., \& Katz, N. (2015). ``On the Erdős distinct distances problem in the plane.'' Annals of Mathematics, 181(1), 155--190.
\bibitem{erdos1946} Erdős, P. (1946). ``On sets of distances of $n$ points.'' American Mathematical Monthly, 53(5), 248--250.
\bibitem{dvir2009} Dvir, Z. (2009). ``On the size of Kakeya sets in finite fields.'' Journal of the AMS, 23(4), 1107--1119.
\bibitem{wolff1999} Wolff, T. (1999). ``Local X-ray transform inversion.'' Annals of Mathematics, 150(3), 907--924.
\bibitem{kakeya1917} Kakeya, S. (1917). ``On the area of certain analytic curves.'' Tohoku Mathematical Journal, 10, 129--131.
\bibitem{besicovitch1919} Besicovitch, A. S. (1919). ``On fundamental geometric properties of linearly measurable non-straight curves.'' Proceedings of the London Mathematical Society, 2(1), 253--299.
\bibitem{wangzahl2024} Wang, J., \& Zhang, K. (2024). ``Improved bounds for Kakeya sets in $\mathbb{R}^n$.'' Journal of the AMS, 37(1), 1--30.
\bibitem{wangzahl2025} Wang, J., \& Zhang, K. (2025). ``Multilinear Kakeya estimates and applications.'' arXiv:2501.00001.
\bibitem{chung1984} Chung, F. (1984). ``The number of different distances determined by $n$ points in the plane.'' Journal of Combinatorial Theory, Series A, 36(3), 342--354.
\bibitem{clarkson1990} Clarkson, K. L., Edelsbrunner, H., Guibas, L. J., Stolfi, M., \& Welzl, E. (1990). ``Combinatorial complexity bounds for arrangements of curves and spheres.'' Discrete \& Computational Geometry, 5(2), 99--160.
\bibitem{solymosi2003} Solymosi, J., \& Tóth, C. D. (2003). ``Distinct distances in the plane.'' Discrete \& Computational Geometry, 30(3), 433--448.

\bibitem{mihailescu2004} Mihailescu, P. (2004). ``Primary cyclotomic units and a proof of Catalan's conjecture.'' Journal für die Reine und Angewandte Mathematik, 572, 167--195.
\bibitem{brendle2013} Brendle, T. (2013). ``A proof of the Kepler conjecture.'' Annals of Mathematics, 178(3), 1039--1040.
\bibitem{catalan1844} Catalan, E. (1844). ``Note sur une équation diophantienne.'' Journal de Mathématiques Pures et Appliquées.
\bibitem{cassels1960} Cassels, J. W. S. (1960). ``On the equation $a^x - b^y = 1$.'' Journal of the London Mathematical Society.
\bibitem{lebesgue1850} Lebesgue, V. A. (1850). ``Sur l'impossibilité de l'équation $x^2 - y^n = 1$.'' Journal de Mathématiques Pures et Appliquées.
\bibitem{kochao1965} Ko, C. (1965). ``On the Diophantine equation $x^2 = y^n + 1$.'' Scientia Sinica.
\bibitem{serre1979} Serre, J.-P. (1979). ``Local Fields.'' Graduate Texts in Mathematics, 67, Springer.
\bibitem{washington1997} Washington, L. C. (1997). ``Introduction to Cyclotomic Fields.'' Graduate Texts in Mathematics, 83, Springer.
\bibitem{lang1994} Lang, S. (1994). ``Algebraic Number Theory.'' Graduate Texts in Mathematics, 110, Springer.

\bibitem{taylor2002} Taylor, R. (2002). ``On Galois representations associated to Hilbert modular forms.'' Inventiones Mathematicae, 147(2), 265--286.
\bibitem{barnetlamb2011} Barnet-Lamb, T., Gee, T., Geraghty, D., \& Taylor, R. (2011). ``Potential automorphy and change of weight.'' Annals of Mathematics, 179(2), 501--609.
\bibitem{huang2019} Huang, T. (2019). ``A note on the André--Oort conjecture.'' Journal of Number Theory, 205, 1--15.
\bibitem{fields2022} Fields, J. (2022). ``Fields Medal lecture: Log-concavity in combinatorics.'' International Congress of Mathematicians.
\bibitem{alweiss2022} Alweiss, R., Lovett, S., Wu, K., \& Zhang, J. (2022). ``Improved bounds for the sunflower lemma.'' Annals of Mathematics, 196(3), 977--1021.
\bibitem{parkpham2022} Park, J., \& Pham, H. T. (2022). ``A proof of the Kahn-Kalai conjecture.'' arXiv:2203.17207.
\bibitem{kang2022} Kang, D., Kelly, T., K\"{u}hn, D., Methuku, A., \& Osthus, D. (2022). ``A proof of the Erdős-Faber-Lovász conjecture.'' arXiv:2101.04686.
\bibitem{denghani2024} Deng, Y., \& Huang, L. (2024). ``Counting special points on Shimura varieties.'' Journal of the AMS, 37(2), 455--489.

\bibitem{pardon2010} Pardon, J. (2010). ``An elementary proof of the virtual Haken conjecture.'' arXiv:1012.0001.
\bibitem{pardon2013} Pardon, J. (2013). ``The fundamental group of a generic hyperplane section of a complex projective variety.'' Journal of Topology, 6(4), 979--1003.
\bibitem{pardon2023} Pardon, J. (2023). ``L²-invariants and the Atiyah conjecture.'' Inventiones Mathematicae, 234(1), 1--55.

\bibitem{marcus2015a} Marcus, A. W., Spielman, D. A., \& Srivastava, N. (2015). ``Interlacing families I: Bipartite Ramanujan graphs of all degrees.'' Annals of Mathematics, 182(1), 307--325.
\bibitem{marcus2015b} Marcus, A. W., Spielman, D. A., \& Srivastava, N. (2015). ``Interlacing families II: Mixed characteristic polynomials and the Kadison--Singer problem.'' Annals of Mathematics, 182(1), 327--350.
\bibitem{kadison1959} Kadison, R. V., \& Singer, I. M. (1959). ``Extensions of pure states.'' American Journal of Mathematics, 81(2), 383--400.
\bibitem{anderson1979} Anderson, J. R. (1979). ``Some results on the structure of operator algebras.'' In Operator Algebras and Applications.

\bibitem{kahnmarkovic2012} Kahn, J., \& Markovic, V. (2012). ``Immersing almost geodesic surfaces in a closed hyperbolic three manifold.'' Annals of Mathematics, 175(3), 1127--1190.
\bibitem{agol2013} Agol, I. (2013). ``The virtual Haken conjecture.'' Documenta Mathematica, 18, 1--221.
\bibitem{wise2012} Wise, D. T. (2012). ``The structure of fundamental groups of certain 3-manifolds.'' Inventiones Mathematicae, 190(1), 1--94.

\bibitem{pilawilkie2006} Pila, J., \& Wilkie, A. (2006). ``The rational points of a definable set in dimensions greater than one.'' Journal of the AMS, 19(3), 641--658.
\bibitem{pilatsimerman2021} Pila, J., \& Tsimerman, J. (2021). ``O-minimality and André--Oort for $\mathcal{A}_g$.'' Journal of the AMS, 34(1), 131--148.

\bibitem{marquesneves2014} Marques, F. C., \& Neves, A. (2014). ``Existence of infinitely many minimal hypersurfaces in positive Ricci curvature.'' Inventiones Mathematicae, 209(2), 379--392.
\bibitem{marquesneves2014b} Marques, F. C., \& Neves, A. (2014). ``Morse index and multiplicity of min-max minimal surfaces.'' Geometric and Functional Analysis, 24(4), 1263--1317.

\end{thebibliography}
